\section{Trabajo autónomo: mantenimiento, instrumentales e integración real}
\label{sec:autonomo}
\textbf{19--20 de octubre de 2026.} Se implementan las actividades que no necesitan
juicios humanos ni una clave del proveedor. Los informes nuevos se conservan en
\code{validation/autonomous/}; los bloques anteriores permanecen como evidencia
histórica. El catálogo del usuario conserva sus 258 canciones: las ingestas de
prueba se realizan en catálogos aislados.

\subsection{E17/P17/P20: cuotas, retención y restauración}
Antes se retenían indefinidamente generaciones, trabajos y medios. Ahora hay
cuotas configurables: catálogo 512 MiB, ingesta 2 GiB y copia 2 GiB por defecto.
La admisión reserva conservadoramente 300 MiB por trabajador activo y por el
siguiente trabajo; la cola mantiene su límite de 32 trabajos. Una publicación
comprueba el uso antes y después de escribir y no cambia el puntero si excede
la cuota. El guard del trabajo comprueba el uso durante el progreso. No se trata
de una cuota impuesta por el sistema operativo: la descarga puede escribir antes
de que un control la detecte; los límites de vídeo y archivo siguen vigentes.
Los pesos de modelos y las copias se administran aparte del catálogo/ingesta.

\code{python -m maintenance cleanup} sólo muestra el plan. La opción
\code{--apply} caduca versiones con más de 30 días, conservando al menos las tres
más recientes y siempre la actual. Conserva medios referenciados por cualquier
versión retenida; trabajos terminales antiguos y archivos huérfanos pueden
retirarse. Los directorios y medios se mueven a cuarentena dentro de su raíz;
el período de recuperación comienza en ese momento. La eliminación permanente
requiere otra operación explícita, \code{purge-trash --apply}, tras la retención.
Los enlaces a versiones caducadas dejan de funcionar: es una contrapartida
del límite de almacenamiento, no una migración transparente de enlaces antiguos.
No se aplicó una eliminación permanente a datos del usuario.

Las rutas se resuelven dentro de raíces administradas y se rechazan enlaces
simbólicos y junctions. Si una generación retenida está corrupta se aborta el
plan, en lugar de suponer que sus medios son prescindibles. El mantenimiento
requiere una cola sin trabajos activos y toma bloqueos en el mismo orden que
publicación. Se caduca el historial antes de mover archivos: una caída puede
dejar archivos sobrantes; no cambia el puntero actual para apuntar a archivos
eliminados. Las dos bases SQLite no constituyen una transacción distribuida.

\code{backup --path} crea un ZIP con catálogo publicado, medios referenciados y
ambas bases SQLite serializadas; excluye claves, configuración privada, pesos,
trabajos temporales y cuarentena. Verifica SHA-256 de cada entrada antes de
publicar el ZIP. \code{restore --path --destination} exige una carpeta nueva,
valida inventario, hashes, integridad SQLite, índices y presencia de medios;
no sobrescribe ni cambia automáticamente el servicio en ejecución. Hay una copia
real del catálogo actual y una restauración verificada en una carpeta aislada.
Para una recuperación se configuran ambas rutas restauradas y se reinicia el
servicio. Los hashes detectan corrupción; no autentican un ZIP de un tercero.

\subsection{E19/P06/P19: instrumentales sin letras ficticias}
El mapa v1 imponía una fila por posición en ambos índices. El mapa v2 conserva
un mapa completo de catálogo y mapas independientes para letras y audio. Se
validan posiciones, IDs únicos, disponibilidad del texto y cobertura acústica.
Las generaciones v1 siguen siendo legibles, sin reconstruir sus vectores.
Una canción sin letra no tiene posición ni vector en el índice lírico; se
mantienen audio y metadatos. Los workers reciben el mapa de su modalidad dentro
del snapshot fijado, evitando atribuir a una canción el vector de otra.

La API admite \code{instrumental:true}; la interfaz ofrece una casilla con esa
decisión. Se omiten subtítulos, Whisper y Nomic para ese trabajo y se registra
que la condición instrumental fue declarada, no verificada acústicamente. Si no
se declara instrumental, una transcripción vacía continúa produciendo error.
La búsqueda acústica puede recuperar el registro; una búsqueda lírica no lo
recupera desde el índice de letras. Un candidato sin texto no recibe aprobación
lírica por el mero hecho de coincidir en audio. La tarjeta explica la ausencia
de letra y no presenta el sonido como auditado.

\subsection{E18/P01/P19: entorno y modelos comprobados realmente}
Se instala \code{.venv-ingestion} separado del entorno histórico. Se resuelven
las versiones completas en \code{requirements/ingestion.lock}; el cierre de
dependencias de modelos se actualiza en \code{requirements/models.lock}.
FFmpeg se aprovisiona desde \code{imageio-ffmpeg} a una ruta local explícita y
no se cambia el PATH global. Whisper recibe muestras decodificadas a 16 kHz;
CLAP usa muestras a 48 kHz. Los caches de pesos quedan en \code{.tools/}.
\code{pip check}, las importaciones y las regresiones se comprueban en el entorno.

La primera inferencia real descubrió una incompatibilidad de Nomic con
Transformers 5: faltaba \code{get_extended_attention_mask}. La solución limita
Transformers a la serie 4 compatible, con 4.57.6 verificada. No se parchea código
descargado ni se oculta la excepción. Se fijan revisiones de Nomic, su código
remoto y CLAP en \code{requirements/model-revisions.json}. Las publicaciones
registran esas revisiones locales; no se atribuyen retroactivamente al corpus
heredado, cuya procedencia exacta sigue sin demostrarse.

La ingesta técnica usa un vídeo de prueba de yt-dlp: descarga real, conversión
MP3, tres ventanas CLAP y publicación aislada con 259 canciones, 258 vectores
líricos y 259 acústicos. Nomic genera un vector normalizado de 256 dimensiones
y Whisper base procesa dos segundos de silencio generado en CPU. El silencio
comprueba ejecución, no exactitud de transcripción de canciones. Los workers
reales ejecutan consultas e índices con muestras iniciales y tres muestras
calientes; éstas pueden usar caché y no equivalen a nuevas inferencias.
Se registra RSS del proceso. No se mide rendimiento del router/juez ni calidad
musical a partir de estas comprobaciones.
El vídeo antiguo estaba retirado; la sustitución se documenta en el
\href{https://github.com/yt-dlp/yt-dlp/issues/12263}{repositorio de yt-dlp}.

\subsection{Interfaz y documento}
La prueba de navegador usa cola, publicación y FFmpeg reales con IA de control,
en un catálogo aislado y con un aviso visible. Comprueba progreso, búsqueda
simultánea, actualización de conteos, recuperación del último trabajo tras
recargar y tarjeta instrumental con reproducción. Se guardan capturas y se
corrige la presentación de la casilla y la conservación de años sin decimales.
La compilación PDF utiliza MiKTeX; se corrigen saltos de línea de rutas técnicas,
se resuelven referencias y se revisan páginas renderizadas. El informe de
compilación distingue esto de las comprobaciones estáticas anteriores.

\subsection{Lo que todavía necesita participación del usuario}
Configurar \code{GROQ_API_KEY} localmente para evaluar router, juez y refinador
reales; no compartir la clave en la conversación. Completar y atribuir las
anotaciones humanas, escuchar las muestras acústicas y revisar las transcripciones
de música antes de declarar relevancia o fidelidad. Después se compara el
refinamiento y se decide si su beneficio compensa coste y latencia. El acceso
a Groq y los juicios humanos no se sustituyen por inferencias de control.
La política de retención es configurable y conservadora; puede ajustarse al uso
real sin que eso impida entregar las herramientas y pruebas de mantenimiento.

\subsection{Compatibilidad del proveedor durante la prueba local}
El 20 de octubre de 2026 se verificó que la credencial permite consultar Groq, pero el modelo Scout configurado devuelve \texttt{model\_not\_found}. Se seleccionó en la configuración local \texttt{openai/gpt-oss-120b}, incluido en los modelos accesibles de la cuenta. También se explicitaron los cuatro nombres de campos obligatorios del planificador y la obligación de conservar las listas vacías: la respuesta anterior empleaba un nombre alternativo incompatible. La consulta ``triste bajo la lluvia produjo un plan válido con estado \texttt{ready}. Esta prueba confirma conexión y contrato, pero no demuestra calidad musical ni valida el flujo completo de búsqueda. Queda pendiente comparar calidad, latencia y coste del modelo con consultas revisadas por personas.
